\begin{table}[H]
\centering
\caption{From Greenwood--Sammon to the Swedish implementation}\label{tab:replication_map}
\small
\setlength{\tabcolsep}{5pt}
\begin{tabular}{p{2.0cm}p{4.65cm}p{6.1cm}}
\toprule
Component & Anchor study & Swedish analysis \\
\midrule
Return pattern & S\&P 500 additions/deletions; announcement, implementation and full windows & OMXSB; same event-window questions, explicitly defined daily-sum CAR \\
Time comparison & Long US history and decade means & Four Swedish calendar blocks and two pooled decades; incomplete 2026 \\
Composition & Period indicators and four firm characteristics & Same regression structure; Swedish measurement proxies and common samples \\
Index demand & Tracking-trade classification; quarterly holdings & Dated mandate groups; available pre/post reports; FI comparison \\
Other index activity & Migrations within S\&P index families & Coincident MSCI changes; a distinct overlap diagnostic \\
Anticipation & Pre-announcement returns and longer starting window & Pre-return windows and AD-20 to ED+1 ratio; historical eligibility not fully reconstructed \\
Liquidity & Effective spreads and implementation costs & Closing quoted spreads and daily volume concentration \\
Extensions & Anchor comparison provides the baseline & Fixed-duration timing, within-firm share classes and separate holdings reporting \\
\bottomrule
\end{tabular}
\begin{minipage}{\linewidth}\footnotesize\vspace{5pt}
The mapping distinguishes conceptual replication from exact numerical reproduction. Detailed source-version references and adaptations appear in Appendix~\ref{app:replication}. Differences in periods, sample selection, holdings dates and return conventions prevent interpreting raw cross-country differences as causal effects.
\end{minipage}
\end{table}
